% Consenso informato (adulti) - progetto brAInc
% Testo identico al template fornito dal corso: compilati solo i campi richiesti.
\documentclass[10pt,a4paper]{article}

\usepackage[T1]{fontenc}
\usepackage[utf8]{inputenc}
% Sillabazione disattivata: il pacchetto italiano di babel non è installato
\hyphenpenalty=10000
\exhyphenpenalty=10000
\emergencystretch=3em
\usepackage{tgtermes}
\usepackage{tgheros}
\usepackage{amssymb}
\usepackage[a4paper,top=2.5cm,bottom=2cm,left=2cm,right=2cm]{geometry}
\usepackage{enumitem}
\usepackage{microtype}
\usepackage{url}
\usepackage{needspace}
\urlstyle{same}

\setlength{\parindent}{0pt}
\setlength{\parskip}{0pt}
\pagestyle{empty}
\raggedbottom

\setlist[itemize]{nosep,leftmargin=2.2em,label=\textbullet}

\newcommand{\progetto}{brAInc}
\newcommand{\corso}{Fondamenti di Human Computer Interaction}
\newcommand{\contatto}{\url{biancaroberta.ianosel@mail.polimi.it}}

\newcommand{\sezione}[1]{\par\needspace{4\baselineskip}\textbf{#1}\par}
\newcommand{\scelta}[1]{%
  \par{\sffamily\itshape\small $\square$ acconsente $\square$ non acconsente\par #1\par}}

\begin{document}

\vspace*{1em}
\textbf{DICHIARAZIONE DI CONSENSO AL TRATTAMENTO DEI DATI PERSONALI AI SENSI\\
DELL'ART. 13 DEL REGOLAMENTO UE N. 679/2016 DEL 27 APRILE 2016}

\medskip
{\large Ringraziandoti per la partecipazione alla ricerca per il progetto didattico
\progetto{} del corso \corso{} del Politecnico di Milano, forniamo di seguito, ai sensi
dell'art.~13 del GDPR (General Data Protection Regulation), regolamento UE n.~679/2016 del
27 aprile 2016, alcune informazioni sul trattamento dei dati che saranno raccolti durante lo
studio.\par}

\medskip
\sezione{Breve descrizione del progetto di ricerca}
Il progetto, svolto dagli studenti Lorenzo Corallo, Bianca Roberta Ianosel e Gabriele Viganò,
studia come le persone leggono e comprendono i termini e condizioni e le privacy policy dei
servizi online. La partecipazione consiste in un'intervista,
registrata in audio, durante la quale potranno essere scattate fotografie.
Le trascrizioni anonimizzate e, per chi vi acconsente, le fotografie saranno pubblicate sul
repository GitHub pubblico del progetto.

\bigskip
\sezione{Base giuridica e natura del conferimento dei dati}
La base giuridica di tale trattamento è da rinvenirsi, ai sensi dell'art.~6, 1° comma,
lett.~a) del Regolamento, nel Suo consenso, libero e facoltativo.\\
La partecipazione al progetto è facoltativa. Il conferimento dei dati richiesti tuttavia è
indispensabile allo svolgimento dello studio: il rifiuto di conferirli non Le consentirà di
parteciparvi.
\sezione{Modalità di trattamento e natura del conferimento dei dati}
Il trattamento dei dati avverrà attraverso strumenti informatici e telematici, con logiche di
organizzazione ed elaborazione correlate esclusivamente alle finalità della ricerca e,
comunque, in modo da garantire la sicurezza e la riservatezza degli stessi in conformità alla
normativa vigente.\\
Le registrazioni audio e video delle interazioni verbali potranno essere trascritte,
conservate, informatizzate e raccolte unicamente ai fini scientifici. In tal caso chi conduce
lo studio si impegna a:
\begin{itemize}
  \item trattare i dati tramite i ricercatori specificatamente designati;
  \item curare le successive fasi di elaborazione e di memorizzazione dei dati raccolti in modo
        da non identificare direttamente gli interessati, procedendo all'anonimizzazione di
        tutte le informazioni;
  \item comunicare i dati e/o memorizzarli, in forma anonima, in un database accessibile a
        tutti i partecipanti al progetto di ricerca;
  \item diffondere le risposte fornite soltanto in modo aggregato e/o anonimo.
\end{itemize}

\medskip
\sezione{Comunicazione e diffusione}
I dati personali, se oggetto di comunicazione e/o diffusione (per esempio per mezzo di
pubblicazioni scientifiche, congressi, seminari, lezioni ecc..), saranno in ogni caso
previamente resi anonimi e trattati in forma aggregata in modo da escludere qualsiasi
possibilità di identificazione del partecipante.\\
Le precisiamo che, nel caso in cui le registrazioni dovessero rivelarsi socialmente
apprezzabili e utilizzabili per presentazioni a eventi di ricerca e conferenze scientifiche,
queste potrebbero essere -- con il Suo consenso - oggetto di comunicazione e/o diffusione.\\
Si precisa che la comunicazione o diffusione dei dati sopra descritti avverranno, con il Suo
consenso, solo previa valutazione della pertinenza e non eccedenza del trattamento rispetto
alle finalità della raccolta o qualora la mancata pubblicazione di quanto emerso nelle
interazioni verbali incida negativamente sulla qualità della ricerca/studio. La registrazione
delle interazioni verbali e/o la trascrizione del contenuto potrebbe essere comunicata ai
partecipanti al progetto ed essere oggetto di diffusione/pubblicazione (es: su riviste
scientifiche, internet, database accessibili ad altri ricercatori).\\
Si precisa, infine, che il Suo consenso alla comunicazione e/o diffusione delle riprese
audio-video delle interazioni verbali implica la concessione di una licenza non esclusiva,
senza limiti di durata e per tutto il mondo, trasferibile a terzi, per l'utilizzazione delle
immagini. Tale licenza include i diritti di cui agli artt.~12 segg.\ della legge n.~633/1941,
compresi a titolo esemplificativo e non esaustivo: diritto di pubblicazione; diritto di
riproduzione in qualunque modo o forma; diritto di trascrizione, montaggio, adattamento,
elaborazione e riduzione; diritto di comunicazione e distribuzione al pubblico, comprendente i
diritti di proiezione, trasmissione e diffusione, anche in versione riassuntiva e/o ridotta,
con qualsiasi mezzo tecnico, diritto di conservare copia delle immagini, anche in forma
elettronica e su qualsiasi supporto tecnologico noto o di futura invenzione per le finalità e
nei limiti sopra definiti. L'uso delle immagini non dà diritto ad alcun compenso. È in ogni
caso esclusa qualunque utilizzazione del ritratto che possa arrecare pregiudizio all'onore,
alla reputazione o al decoro della persona ritratta, ripresa o registrata.

\bigskip
\sezione{Periodo di conservazione dei dati}
I dati raccolti saranno utilizzati con strumenti informatici e telematici e saranno conservati
esclusivamente per il tempo strettamente necessario per la trascrizione dei contenuti e per la
loro anonimizzazione e, comunque, non oltre 2 anni.

\bigskip
\sezione{Diritti degli interessati}
In qualità di soggetto interessato può chiedere in qualsiasi momento al Titolare:
\begin{itemize}
  \item di ritirarsi dallo studio in qualsiasi momento
  \item la conferma dell'esistenza o meno di dati personali che lo riguardano;
  \item l'accesso ai suoi dati personali e alle informazioni relative agli stessi;
  \item la rettifica dei dati inesatti o l'integrazione di quelli incompleti;
  \item la cancellazione dei dati personali che la riguardano (al verificarsi di una delle
        condizioni indicate nell'art.~17, paragrafo 1 del Regolamento e nel rispetto delle
        eccezioni previste nel paragrafo 3 dello stesso articolo);
  \item la limitazione del trattamento dei suoi dati personali (al ricorrere di una delle
        ipotesi indicate nell'art.~18, paragrafo 1 del Regolamento);
  \item la trasformazione in forma anonima o il blocco dei dati trattati in violazione di
        legge, compresi quelli di cui non è necessaria la conservazione in relazione agli scopi
        per i quali i dati sono stati raccolti o successivamente trattati
\end{itemize}

\medskip
In qualità di soggetto interessato ha inoltre diritto di opporsi in tutto o in parte:
\begin{itemize}
  \item per motivi legittimi al trattamento dei dati personali che lo riguardano, ancorché
        pertinenti allo scopo della raccolta;
  \item al trattamento di dati personali che lo riguardano a fini di invio di promozione di
        iniziative formative ed eventi culturali.
\end{itemize}

\medskip
Tali diritti sono esercitabili rivolgendosi a \url{privacy@polimi.it}.\\
Qualora ritenga che i suoi diritti siano stati violati dal titolare e/o da un terzo, ha il
diritto di proporre reclamo all'Autorità per la protezione dei dati personali e/o ad altra
autorità di controllo competente in forza del Regolamento.

\bigskip\bigskip
\sezione{CONSENSO AL TRATTAMENTO DEI DATI PERSONALI}
Il/La sottoscritto/a \dotfill\\
nato/a a \makebox[3.5cm]{\dotfill} il\makebox[4.5cm]{\dotfill} residente a
\makebox[5cm]{\dotfill}\\
via\makebox[9cm]{\dotfill} n\makebox[1.2cm]{\dotfill} ;\\
presa visione dell'informativa ai sensi dell'art. dell'art.~13 del GDPR (General Data Protection
Regulation), regolamento UE n.~679/2016 del 27 aprile 2016\\
\textbf{AUTORIZZA}\\
Gli studenti responsabili dello studio (Lorenzo Corallo, Bianca Roberta Ianosel, Gabriele
Viganò) al trattamento dei dati personali del partecipante per le finalità indicate
nell'informativa di cui sopra.

\scelta{al trattamento dei propri dati, anche di natura particolare, per finalità di ricerca
statistica e scientifica con le modalità e per gli scopi descritti [adesione al progetto];}
\scelta{al trattamento dei propri dati, anche di natura particolare, per finalità di ricerca
statistica e scientifica, incluse la comunicazione e la diffusione delle interviste, con le
modalità e per gli scopi descritti [diffusione delle interviste prive di dati
identificativi];}
\scelta{l'uso, la riproduzione e la pubblicazione con ogni con ogni mezzo tecnico delle immagini
riprese e la possibilità che le registrazioni possano essere oggetto di pubblica
rappresentazione attraverso manifestazioni, pubblicazione su web e social network,
downloading, diffusione su supporto ottico o magnetico.}

\bigskip
Luogo e data \makebox[6cm]{\dotfill}

\bigskip
(firma leggibile) \makebox[6cm]{\dotfill}

\end{document}
